\documentclass[12pt,a4paper]{article}
\usepackage[utf8]{inputenc}
\usepackage[portuges]{babel}
\usepackage[T1]{fontenc}
\usepackage{amsmath}
\usepackage{amsfonts}
\usepackage{amssymb}
\usepackage{graphicx}
\usepackage{float}
\usepackage{hyperref}
\usepackage{geometry}
\usepackage{multirow}
\usepackage{array}
\usepackage{booktabs}
\usepackage{xcolor}
\usepackage{listings}
\usepackage{caption}

\geometry{left=3cm, right=2cm, top=3cm, bottom=2cm}

\hypersetup{
    colorlinks=true,
    linkcolor=blue,
    filecolor=magenta,      
    urlcolor=cyan,
    pdftitle={Relatório - Algoritmos de Busca},
    pdfpagemode=FullScreen,
}

\lstset{
    language=Python,
    basicstyle=\ttfamily\small,
    keywordstyle=\color{blue},
    commentstyle=\color{green!60!black},
    stringstyle=\color{orange},
    numbers=left,
    numberstyle=\tiny\color{gray},
    stepnumber=1,
    numbersep=5pt,
    backgroundcolor=\color{gray!10},
    showspaces=false,
    showstringspaces=false,
    showtabs=false,
    frame=single,
    tabsize=4,
    captionpos=b,
    breaklines=true,
    breakatwhitespace=false,
    escapeinside={\%*}{*)}
}

\title{%
    \Large UNIVERSIDADE FEDERAL DE PELOTAS – UFPEL \\
    CENTRO DE DESENVOLVIMENTO TECNOLÓGICO (CDTec) \\
    CURSOS DE CIÊNCIA DA COMPUTAÇÃO E ENGENHARIA DE COMPUTAÇÃO \\
    DISCIPLINA DE FUNDAMENTOS DE INTELIGÊNCIA ARTIFICIAL \\
    PROF.: Dr. ANDERSON PRIEBE FERRUGEM
    \vspace{2cm}
    \Huge\textbf{RELATÓRIO DE TRABALHO PRÁTICO} \\
    \Large\textbf{ALGORITMOS DE BUSCA EM ESPAÇO DE ESTADOS} \\
    \Large\textbf{SISTEMA DE NAVEGAÇÃO PARA ENTREGAS URBANAS}
}

\date{}

\begin{document}

\maketitle
\thispagestyle{empty}

\newpage

\section*{Dados do Trabalho}
\begin{tabular}{|p{4cm}|p{10cm}|}
\hline
\textbf{Título} & Sistema de Navegação para Entregas Urbanas \\
\hline
\textbf{Disciplina} & Fundamentos de Inteligência Artificial \\
\hline
\textbf{Professor} & Dr. Anderson Priebe Ferrugem \\
\hline
\textbf{Data de Entrega} & \_\_\_\_/\_\_\_\_/\_\_\_\_\_\_ \\
\hline
\end{tabular}

\vspace{1cm}

\section*{Identificação do Grupo}
\begin{tabular}{|p{4cm}|p{10cm}|}
\hline
\textbf{Nome do Grupo} &  \\
\hline
\textbf{Componentes} & Nome completo – Matrícula \\
 & Nome completo – Matrícula \\
 & Nome completo – Matrícula \\
 & Nome completo – Matrícula \\
 & Nome completo – Matrícula \\
\hline
\end{tabular}

\newpage

\tableofcontents
\newpage

\section{INTRODUÇÃO}

[Contextualização do problema – máximo 20 linhas]

\begin{itemize}
    \item Apresente o problema de roteirização urbana e sua importância para empresas de entrega.
    \item Explique o que são algoritmos de busca em espaço de estados e por que são adequados para este problema.
    \item Descreva brevemente o objetivo do trabalho: implementar e comparar diferentes algoritmos de busca para encontrar a melhor rota entre dois pontos no centro de Pelotas.
    \item Cite as ferramentas utilizadas (Python, NetworkX, etc.).
\end{itemize}

\section{FUNDAMENTAÇÃO TEÓRICA}

[Explique cada um dos conceitos abaixo com suas próprias palavras – máximo 30 linhas]

\subsection{Busca em Espaço de Estados}
\begin{itemize}
    \item Definição de espaço de estados, nós, arestas, custos.
    \item Diferença entre busca cega e busca informada.
\end{itemize}

\subsection{Algoritmos Implementados}
Para cada algoritmo, descreva:

\subsubsection{BFS (Busca em Largura)}
Funcionamento, estrutura de dados utilizada, completude, optimalidade (quando se aplica), complexidades.

\subsubsection{DFS (Busca em Profundidade)}
Funcionamento, estrutura de dados utilizada, completude, optimalidade, complexidades.

\subsubsection{UCS (Busca de Custo Uniforme)}
Funcionamento, estrutura de dados utilizada, completude, optimalidade, complexidades.

\subsubsection{Busca Gulosa (Best-First)}
Funcionamento, papel da heurística, completude, optimalidade.

\subsubsection{A* (A-Estrela)}
Funcionamento, função \( f(n) = g(n) + h(n) \), condições para otimalidade (heurística admissível e consistente).

\subsection{Heurística}
\begin{itemize}
    \item Definição de heurística.
    \item Explique a heurística utilizada (distância em linha reta) e justifique por que ela é admissível para este problema.
\end{itemize}

\section{METODOLOGIA}

[Descreva passo a passo como o trabalho foi desenvolvido – máximo 30 linhas]

\subsection{Modelagem do Grafo}
\begin{itemize}
    \item Como os dados do mapa foram carregados (formato CSV).
    \item Nós e arestas: o que cada nó representa, quais informações foram armazenadas (coordenadas, nomes).
    \item Custos: distância (metros) e tempo (minutos).
\end{itemize}

\subsection{Implementação dos Algoritmos}
\begin{itemize}
    \item Explique como cada algoritmo foi implementado (estruturas de dados utilizadas: pilha, fila, heapq).
    \item Detalhe como o registro de nós expandidos e fronteira máxima foi realizado.
    \item Como a detecção de ciclos foi tratada.
\end{itemize}

\subsection{Cenários de Teste}
\begin{itemize}
    \item \textbf{Cenário A (Minimizar Distância):} origem, destino, critério utilizado.
    \item \textbf{Cenário B (Minimizar Tempo):} origem, destino, critério utilizado.
    \item Justifique a escolha dos pontos de origem e destino.
\end{itemize}

\subsection{Visualização}
Explique como os caminhos foram visualizados (biblioteca matplotlib, função implementada).

\section{RESULTADOS}

\subsection{Tabela Comparativa – Cenário A (Distância)}

\begin{table}[H]
\centering
\begin{tabular}{|l|c|c|c|c|}
\hline
\textbf{Algoritmo} & \textbf{Custo total (m)} & \textbf{Nós expandidos} & \textbf{Fronteira máxima} & \textbf{Solução ótima?} \\
\hline
BFS &  &  &  &  \\
\hline
DFS &  &  &  &  \\
\hline
UCS &  &  &  &  \\
\hline
Gulosa &  &  &  &  \\
\hline
A* &  &  &  &  \\
\hline
\end{tabular}
\caption{Resultados dos algoritmos para o Cenário A (minimizar distância)}
\end{table}

\subsection{Tabela Comparativa – Cenário B (Tempo)}

\begin{table}[H]
\centering
\begin{tabular}{|l|c|c|c|c|}
\hline
\textbf{Algoritmo} & \textbf{Custo total (min)} & \textbf{Nós expandidos} & \textbf{Fronteira máxima} & \textbf{Solução ótima?} \\
\hline
BFS &  &  &  &  \\
\hline
DFS &  &  &  &  \\
\hline
UCS &  &  &  &  \\
\hline
Gulosa &  &  &  &  \\
\hline
A* &  &  &  &  \\
\hline
\end{tabular}
\caption{Resultados dos algoritmos para o Cenário B (minimizar tempo)}
\end{table}

\subsection{Visualizações}

[Insira as imagens geradas para cada algoritmo em cada cenário, com legendas explicativas]

\begin{figure}[H]
\centering
% \includegraphics[width=0.8\textwidth]{figura1.png}
\caption{Caminho encontrado pelo BFS (Cenário A)}
\end{figure}

\begin{figure}[H]
\centering
% \includegraphics[width=0.8\textwidth]{figura2.png}
\caption{Caminho encontrado pelo DFS (Cenário A)}
\end{figure}

\begin{figure}[H]
\centering
% \includegraphics[width=0.8\textwidth]{figura3.png}
\caption{Caminho encontrado pelo UCS (Cenário A)}
\end{figure}

\begin{figure}[H]
\centering
% \includegraphics[width=0.8\textwidth]{figura4.png}
\caption{Caminho encontrado pela Busca Gulosa (Cenário A)}
\end{figure}

\begin{figure}[H]
\centering
% \includegraphics[width=0.8\textwidth]{figura5.png}
\caption{Caminho encontrado pelo A* (Cenário A)}
\end{figure}

\section{DISCUSSÃO}

[Responda às perguntas abaixo de forma argumentada, com base nos resultados obtidos – máximo 50 linhas]

\subsection{Completude}
Todos os algoritmos encontraram solução? Se algum falhou, explique por quê.

\subsection{Optimalidade}
Quais algoritmos encontraram a solução de menor custo em cada cenário? Por que BFS encontrou o caminho de menor distância mesmo sem considerar custos? (Relacione com o fato de todas as arestas terem custos iguais ou diferentes.)

\subsection{Complexidade de Tempo}
Qual algoritmo expandiu mais nós? Qual expandiu menos? Relacione com a complexidade teórica de cada um e com as características do grafo.

\subsection{Complexidade de Espaço}
Qual algoritmo manteve a maior fronteira (maior consumo de memória)? Explique por que isso acontece (DFS tem baixa fronteira, BFS tem alta, etc.).

\subsection{Impacto da Heurística}
Compare os resultados da Busca Gulosa e do A*. A heurística utilizada é admissível? Isso influenciou os resultados? A Busca Gulosa encontrou a solução ótima? Se não, por quê?

\subsection{Análise Prática}
Se você fosse implementar um sistema de navegação em tempo real para entregadores, qual algoritmo escolheria? Considere o equilíbrio entre optimalidade, velocidade (nós expandidos) e consumo de memória. Justifique.

\section{SISTEMA DE ROTEIRIZAÇÃO DESENVOLVIDO}

\subsection{Descrição do Sistema}
\begin{itemize}
    \item Explique o funcionamento do script \texttt{roteirizador.py}.
    \item Quais são as funcionalidades disponíveis? (menu interativo, escolha de origem/destino, critério, algoritmo, visualização, salvamento de log).
\end{itemize}

\subsection{Exemplos de Execução}
[Insira prints da execução do sistema com diferentes configurações]

\subsubsection{Exemplo 1:} Origem CD\_CENTRO → Destino PRACA\_PEDRO\_OSORIO, critério distância, algoritmo A*

[Print da tela + imagem do mapa gerado]

\subsubsection{Exemplo 2:} Origem CD\_CENTRO → Destino PRACA\_PEDRO\_OSORIO, critério tempo, algoritmo UCS

[Print da tela + imagem do mapa gerado]

\subsubsection{Exemplo 3:} (Outra combinação relevante)

[Print da tela + imagem do mapa gerado]

\subsection{Código-Fonte (Trechos Relevantes)}
[Insira os trechos mais importantes do código com explicações]

\begin{lstlisting}[caption=Implementação do algoritmo A*]
def busca_a_star(grafo, inicio, objetivo, heuristica, custo='distance'):
    """
    Implementação do algoritmo A* para busca de caminhos
    Retorna: (caminho, custo_total, nos_expandidos, fronteira_max)
    """
    fronteira = []  # Fila de prioridade
    heapq.heappush(fronteira, (0 + heuristica(inicio, objetivo), 0, inicio, [inicio]))
    
    explorados = set()
    nos_expandidos = 0
    fronteira_max = 0
    
    while fronteira:
        fronteira_max = max(fronteira_max, len(fronteira))
        f, g, no_atual, caminho = heapq.heappop(fronteira)
        
        if no_atual in explorados:
            continue
            
        nos_expandidos += 1
        explorados.add(no_atual)
        
        if no_atual == objetivo:
            return caminho, g, nos_expandidos, fronteira_max
        
        for vizinho in grafo.neighbors(no_atual):
            if vizinho not in explorados:
                custo_aresta = grafo[no_atual][vizinho][custo]
                novo_g = g + custo_aresta
                novo_f = novo_g + heuristica(vizinho, objetivo)
                novo_caminho = caminho + [vizinho]
                heapq.heappush(fronteira, (novo_f, novo_g, vizinho, novo_caminho))
    
    return None, float('inf'), nos_expandidos, fronteira_max
\end{lstlisting}

\section{CONCLUSÃO}

[Síntese dos principais achados – máximo 20 linhas]

\begin{itemize}
    \item Retome brevemente o objetivo do trabalho.
    \item Destaque os principais resultados e conclusões.
    \item Aponte dificuldades encontradas durante o desenvolvimento.
    \item Sugira melhorias futuras (ex.: considerar trânsito em tempo real, múltiplos destinos, restrições de horário, etc.).
\end{itemize}

\section{REFERÊNCIAS}

[Liste as referências utilizadas, incluindo livros, artigos, documentações e sites]

\begin{itemize}
    \item RUSSELL, Stuart; NORVIG, Peter. Inteligência Artificial. 3. ed. Rio de Janeiro: Elsevier, 2013.
    \item Documentação NetworkX. Disponível em: \url{https://networkx.org/}. Acesso em: [data].
    \item Documentação Matplotlib. Disponível em: \url{https://matplotlib.org/}. Acesso em: [data].
    \item AISpace. Tutorial interativo de busca. Disponível em: \url{http://www.aispace.org/search/}. Acesso em: [data].
\end{itemize}

\section{APÊNDICES}

\subsection{Código-Fonte Completo}
[Opcional – pode ser incluído como link para o repositório ou arquivo anexo]

\subsection{Link do Vídeo de Apresentação}
[Inserir link do YouTube não listado ou Google Drive]

\section{DECLARAÇÃO DE AUTORIA}

Declaramos que este trabalho é de nossa autoria e não foram utilizadas fontes não citadas ou permitidas, incluindo código gerado por inteligência artificial sem devida referência.

\vspace{1cm}

\begin{tabular}{|p{6cm}|p{6cm}|}
\hline
\textbf{Nome} & \textbf{Assinatura} \\
\hline
 &  \\
\hline
 &  \\
\hline
 &  \\
\hline
 &  \\
\hline
 &  \\
\hline
\end{tabular}

\vspace{1cm}

Pelotas, \_\_\_\_ de \_\_\_\_\_\_\_\_\_\_\_\_\_ de \_\_\_\_\_\_.

\newpage

\section*{ANEXO – CHECKLIST DE AVALIAÇÃO (PARA USO DO GRUPO)}

Antes de entregar, verifique se:

\begin{table}[H]
\centering
\begin{tabular}{|p{12cm}|c|}
\hline
\textbf{Item} & \textbf{OK?} \\
\hline
Todos os algoritmos foram implementados (BFS, DFS, UCS, Gulosa, A*) & ☐ \\
\hline
Os experimentos foram realizados para os dois cenários (distância e tempo) & ☐ \\
\hline
As tabelas de resultados estão preenchidas & ☐ \\
\hline
As imagens de visualização foram geradas e incluídas & ☐ \\
\hline
As perguntas da seção de discussão foram respondidas & ☐ \\
\hline
O script \texttt{roteirizador.py} está funcionando e foi testado & ☐ \\
\hline
O link do vídeo está funcionando e o vídeo tem entre 12-15 minutos & ☐ \\
\hline
Todos os membros participaram da apresentação & ☐ \\
\hline
O relatório está formatado corretamente (PDF) & ☐ \\
\hline
O arquivo .zip com todos os materiais foi gerado & ☐ \\
\hline
\end{tabular}
\end{table}

\end{document}